\title{LongRoPE: Extending LLM Context Window Beyond 2 Million Tokens}

\begin{abstract}
A substantial context window is a sought-after capability in large language models (LLMs). Nevertheless, expanding context length is hindered by high fine-tuning expenses, the limited availability of lengthy textual data, and catastrophic behaviors arising from unfamiliar token positions. As a result, present-day models with extended context capacities typically cap at approximately 128k tokens.
\end{abstract}

This work presents LongRoPE, a method that, for the first time, pushes the context window of existing pre-trained LLMs up to a remarkable \textbf{2048k} tokens. This is accomplished with as few as 1k fine-tuning steps conducted at 256k training lengths, while preserving the model’s effectiveness on the initial short context window. LongRoPE achieves this through three principal technical advancements: (i) it uncovers and utilizes two distinct types of non-uniformity in positional interpolation via an efficient search, yielding superior fine-tuning initialization and allowing for up to 8$\times$ window extension without fine-tuning; (ii) it adopts a staged approach whereby a model is first fine-tuned at 256k context length, followed by a second round of positional interpolation on this extended model, resulting in a full 2048k window; (iii) it recalibrates LongRoPE specifically at the 8k length to ensure restoration of short-context performance. Comprehensive empirical results using LLaMA2 and Mistral on multiple benchmarks confirm the practical merits of our approach. The models adapted with LongRoPE largely preserve their original architecture, requiring only slight adjustments to positional embeddings, and can leverage nearly all established optimization techniques. Code will be released at \texttt{https://github.com/microsoft/LongRoPE}

\section{Introduction}

While Large Language Models (LLMs) have demonstrated outstanding performance across a range of tasks~\cite{gpt4,llama2}, they are constrained by relatively small context windows, such as the 4096 token limitation in LLaMA2~\cite{llama2}. When text sequences extend beyond this context window, model performance typically deteriorates because the model has not been trained to attend to these extended positions. This shortcoming becomes particularly problematic in key applications, including in-context learning with extensive example sets~\cite{Huang2023FewerIM}, as well as in deploying LLM agents~\cite{park2023generative,madaan2023selfrefine}.

Recent studies have demonstrated that the context window of a pre-trained LLM can be expanded up to approximately 128k through fine-tuning with longer sequences~\cite{longlora,pi,yarn,zhang2024soaring,liu2023scaling}. However, three principal challenges inhibit further scaling of the context window. Firstly, newly introduced and untrained position indices can yield numerous extreme values, which results in distributional shifts and impedes the convergence of fine-tuning processes~\cite{pi}. This complication is intensified in cases where context windows are augmented from 4k to over 1000k, as over 90\% of the position indices are novel. Secondly, fine-tuning typically necessitates training data with text lengths that match the intended extension. Currently, datasets containing exceedingly long samples—specifically, those beyond 1000k tokens—are rare. In addition, the computational demands for training with such lengthy inputs are considerable, requiring extensive time and substantial GPU resources. Thirdly, upon scaling to much larger context windows, attention becomes increasingly diluted, as it must cover vastly more token positions, thereby reducing the model’s effectiveness for original, shorter input lengths~\cite{pi}.

A common strategy to address the initial challenge involves interpolating the RoPE positional embeddings~\cite{rope,pi}, which effectively scales new position indices to fit within the pretrained range, as illustrated in Fig.\ref{fig:overview}. Position Interpolation (PI)~\cite{pi} performs linear interpolation on RoPE's rotary angles using the extension ratio. NTK~\cite{ntk,dynamicntk} proposes performing non-uniform interpolation and extrapolation across the various RoPE dimensions. YaRN~\cite{yarn} separates the RoPE dimensions based on frequency and applies different extrapolation methods—specifically, NTK and linear interpolation—to each corresponding category. Nevertheless, the positional embedding mechanism inherently demonstrates \emph{complex, non-uniform information entropy} within the Transformer framework. Existing methods tend to overlook these subtle irregularities, resulting in unutilized information and consequently constraining the achievable context window.

Section~\ref{sec:analysis} presents two main empirical observations: \textbf{(1)} Successful positional interpolation must account for both variable RoPE dimension scaling and distinct token position effects. In practice, reduced RoPE dimensions and early token locations require minimal interpolation; however, determining the best interpolation approach is contingent on the intended sequence extension length. \textbf{(2)} By incorporating these sources of non-uniformity into the interpolation strategy, it becomes possible to preserve the intrinsic information contained within the original RoPE—particularly across critical dimensions and at essential token positions. This strategy minimizes the information degradation associated with interpolation and offers a superior starting point for subsequent fine-tuning. Additionally, such an approach facilitates up to 8$\times$ longer context windows even in the absence of fine-tuning.

Building upon these results, we present \textbf{LongRoPE}, a robust approach that broadens the context window of LLMs to accommodate more than 2 \emph{million} tokens. The LongRoPE framework is grounded in three main innovations. The first is its comprehensive utilization of multidimensional positional interpolation non-uniformities. LongRoPE determines suitable rescale factors for the rotation angles of RoPE within each dimensional axis, guided by the relative positions of tokens. Given that the identification of rescale factors involves a search space that increases exponentially with the desired extension ratio, LongRoPE adopts an evolutionary search strategy enhanced with two optimization mechanisms to streamline search performance. An illustration of the resulting rescaled RoPE obtained via this method is depicted in Fig.~\ref{fig:overview}.

Subsequently, {LongRoPE} utilizes a resource-efficient and incremental extension approach to reach a 2048k context window, circumventing the necessity for direct fine-tuning on exceedingly lengthy sequences, which are rarely accessible. The method initiates by identifying a 256k sequence length with the pre-trained LLM, followed by fine-tuning at this particular length. Owing to our non-uniform positional interpolation technique—which enables up to an 8$\times$ context length expansion without requiring further fine-tuning—we then perform a secondary search to determine optimal RoPE rescale factors using the already fine-tuned extended LLM. As a result, this pipeline allows LLaMA2 and Mistral~\cite{mistral} to handle context windows as large as 2048k.

In order to prevent loss of effectiveness when reverting to the initial, smaller context window, {LongRoPE} maintains ongoing refinement of RoPE rescale factors in the expanded LLM model. In analogy to the process used for scaling up from 256k to 2048k, we now reduce the context windows to 4k and 8k on the model fine-tuned for 256k, employing our search method to favor minimal positional interpolation. When generating outputs, if the input sequence is under 8k in length, RoPE is modified using the identified rescale parameters.

Our method has been rigorously evaluated through numerous experiments on a range of LLMs and long-context tasks, evidencing its strong performance. The results indicate that LongRoPE consistently sustains low perplexity when evaluated with lengths spanning from 4k up to 2048k, surpassing 90\% accuracy in passkey retrieval, and achieving similar performance to existing models on established benchmarks constrained to a 4096-token context window. Notably, LongRoPE is compatible with any LLM utilizing RoPE embedding. The codebase and LongRoPE-2048k models will be made publicly available.

\section{Non-uniformity in Positional Interpolation}
\label{sec:analysis}

\subsection{Preliminary}

Transformer models require explicit positional information, often in the form of position embedding, to represent the order of input tokens. Our work focuses on the RoPE~\cite{rope} position embedding, which is widely used in recent  LLMs. For a token at position index $n$, its corresponding RoPE encoding can be simplified as follows:

\begin{equation}
		\fontsize{7.8}{7.8} \selectfont
	\label{eq:rope}
	\left[ cos(n\theta_0), sin(n\theta_0), cos(n\theta_1), \cdot\cdot\cdot, cos(n\theta_{d/2-1}), sin(n\theta_{d/2-1}) \right]   
\end{equation}

Here, $d$ refers to the size of the embedding space, $n\theta_i$ indicates the rotary position angle assigned to a token located at position $n$, and $\theta_i=\theta^{-2i/d}$ denotes the frequency for positional rotations. RoPE typically adopts 10000 as the base value for $\theta$.

\textbf{\emph{Context window extension ratio $s$} and positional interpolation}. The extension ratio $s$ is specified as the quotient of the new context window length $L'$ by the original length $L$, i.e., $s=\frac{L'}{L}$.

To extend the context window from $L$ to $L'$, current positional interpolation methods suggest downscaling rotation frequencies $\theta_i$ based on extension ratio $s$. Let $\beta=\theta^{2/d}$, and $\lambda$ denote the actual rescale factor related to $s$, we   unify these positional interpolation methods as follows:

\begin{equation}	
	\fontsize{7.5}{7.5} \selectfont
	\label{eq:unified_rope}
	\begin{aligned}
		\left[cos\left(\frac{n}{\lambda(\beta)^0}\right), sin \left(\frac{n}{\lambda(\beta)^0}\right), cos\left(\frac{n}{\lambda(\beta)^1}\right), \cdot\cdot\cdot,  sin \left(\frac{n}{\lambda(\beta)^{d/2-1}}\right) \right]   
	\end{aligned}
\end{equation}

\textbf{Linear positional interpolation (PI)}. In the PI method~\cite{pi}, the position indices are interpolated in a linear fashion, staying within the original pre-training length boundaries. To achieve an extension by a factor of $s$, the rotation angles corresponding to each position are uniformly scaled down by a multiplier $\lambda=s$ in every RoPE dimension. Nonetheless, such uniform compression causes position encodings to cluster tightly together, which reduces the model's capacity to differentiate between tokens that are near each other. Consequently, PI's performance often degrades when applied at larger extension ratios.

\textbf{NTK-based interpolation and extrapolation}. The works in~\cite{ntk,dynamicntk} approach RoPE from the standpoint of information encoding and utilize Neural Tangent Kernel (NTK) theory~\cite{jacot2018neural,tancik2020fourier}. To address the issue of positional overcrowding in PI, they propose distributing the interpolation load among the different RoPE dimensions. Their method involves scaling lower (high-frequency) dimensions less and higher (low-frequency) dimensions more, which facilitates both interpolation and extrapolation over positions, with $\lambda=s^{i}$. In the updated dynamic NTK model~\cite{dynamicntk}, the extension ratio is dynamically adapted for each position according to the current sequence length. In contrast to PI—which requires fine-tuning—NTK-aware approaches are capable of context window extension without additional fine-tuning, although this is typically constrained to a maximum extension ratio of 4$\times$.

\textbf{YaRN}~\cite{yarn} offers notable advancements in positional interpolation by categorizing RoPE dimensions according to their associated frequencies, applying distinct interpolation approaches to each group. Specifically, dimensions with high frequencies are subjected to extrapolation ($\lambda$=1), low-frequency dimensions utilize linear interpolation (PI), and those within the intermediate frequency range apply the NTK method. A central aspect of YaRN’s methodology is its partitioning of RoPE dimensions; however, this partitioning is currently determined through manual, empirical tuning. As a consequence, this method may not deliver optimal results when generalized to novel LLMs.

\subsection{Study on Non-uniform Positional Interpolation}

Building upon ideas from NTK and YaRN, it is apparent that their improvements stem from introducing non-linear behavior, particularly by distinguishing among frequencies within RoPE dimensions to tailor both interpolation and extrapolation strategies. Nevertheless, existing non-linear components predominantly depend on hand-crafted heuristics. This observation leads us to two fundamental questions: (1) Does the positional interpolation currently in use represent the best possible approach? (2) Could there exist additional, yet unidentified, forms of non-linearity?

In order to address these questions, we apply evolution search (refer to Sec.~\ref{sec:method}) to identify improved non-uniform positional interpolations tailored for LLaMA2-7B. Our search process is steered by perplexity metrics, utilizing five randomly chosen examples from the PG19~\cite{pg19} validation data. Our empirical investigations yield several significant insights.

\textbf{Finding 1}: \textit{RoPE dimensions exhibit substantial non-uniformities, which are not effectively handled by current positional interpolation methods.}

We determine the optimal $\lambda$ for each RoPE dimension as outlined in Eq.~\ref{eq:unified_rope}. Table~\ref{tbl:exp1} reports the perplexity of LLaMA2-7B across various approaches on the PG19 and Proof-pile~\cite{proof-pile} test datasets, prior to any fine-tuning. Our optimization approach achieves considerable gains, indicating that both the conventional linear (PI) and the non-uniform methods (Dynamic-NTK and YaRN) fall short of optimal performance. Interestingly, on PG19, YaRN lags behind both PI and NTK, as it fails to support the intended context window length without fine-tuning the LLM. Specifically, in an 8k context, YaRN’s perplexity markedly increases once the context surpasses 7k tokens.

In our exploration, we found that the rescaling coefficients $\lambda$ in Eq.~\ref{eq:unified_rope} are no longer uniform, contrasting with the constant scaling factor $s$ used in PI, the formulation in NTK, and the group-level approach adopted by YaRN. Employing these non-uniform scaling factors leads to marked enhancements in the language modeling capability of LLaMA2 (as measured by perplexity) for extended context lengths of 8k and 16k tokens, all achieved without additional fine-tuning. The performance gains can be attributed to the positional embeddings' improved retention of the original RoPE structure—particularly within critical dimensions—which in turn makes it easier for the model to distinguish between closely positioned tokens.

\textbf{Finding 2}: \textit{RoPE for the initial tokens in the input sequence should be extrapolated with less interpolation.} 

For the first $\hat{n}$ tokens of input sequences, we propose that their RoPE encoding should involve minimal interpolation. The rationale behind this is that these tokens, owing to their high attention scores, play a pivotal role within the attention layers—an effect similarly reported in Streaming LLM~\cite{streamingllm} and LM-Infinite~\cite{han2023lminfinite}. To examine this hypothesis, we expand the context window to 8k and 16k tokens using PI and NTK methodologies, deliberately excluding the first $\hat n$ tokens ($0, 2, ..., 256$) from interpolation. When $\hat n=0$, the schemes default to the original PI and NTK configurations. Table~\ref{tbl:tokenposition} yields two primary insights: \textbf{(1)} preserving positional encoding without interpolation for the initial tokens tangibly enhances model performance; \textbf{(2)} the ideal count of leading tokens $\hat n$ varies according to the specific context window extension desired.

\textbf{Finding 3}: \textit{Non-uniform positional interpolation effectively extends LLM context window in both fine-tuning and non-fine-tuning settings.} 

Although our non-uniform position interpolation approach demonstrates considerable gains in extension capability at 8k and 16k context lengths even without fine-tuning, pushing to longer context sizes necessitates model fine-tuning. Consequently, we apply fine-tuning to LLaMA2-7B using our optimized RoPE design for a 64k context window (detailed configurations are provided in the Appendix). According to Table~\ref{tbl:finetune}, this method leads to substantial improvements over PI and YaRN, both prior to and following the fine-tuning of LLaMA2-7B. These advantages can be attributed to the more effective non-uniform positional interpolation, which curbs information degradation and delivers superior initialization for fine-tuning procedures.

\textbf{Summary}. Our findings identify two forms of non-uniformity: changes in RoPE dimensions and token positions. Strategically leveraging these non-uniformities within positional interpolation drives marked enhancements in the context extension performance of LLMs.

\section{LongRoPE}

\label{sec:method}
Building on these insights, we propose LongRoPE, which initially employs an effective search strategy designed to harness both types of non-uniformity, subsequently applying this method to increase the LLM context window to over 2 million tokens.

\subsection{Problem Formulation}

These two forms of non-uniformity greatly expand the set of possible solutions and complicate the optimization process. To mitigate these challenges, we recast the multidimensional non-uniform position interpolation optimization as a search problem.

For a LLM targeting a  context window size of $L'$ and lengthy input documents $\mathbf{X}$, where each $\mathbf{x}\in\mathbf{X}$ surpasses $L'$ in token length, we denote the original rotary angle of the $i^{th}$ dimension in RoPE embedding at token position $n$ as  $\frac{n}{\beta^i}$. The optimization problem is then formulated as follows:

\begin{equation}
\begin{gathered}
		\label{eq:problem}

\underset {\mathbf{x}\in\mathbf{X}; \, |\mathbf{x}|\ge L'}{ \operatorname {arg\,min} } \,  \mathcal{L}\, \left( \text{LLM} (\text{RoPE}, \mathbf{X})\right), \text{where }
\\
\scalebox{0.93}{
$\underset {\begin{subarray}{l} i=0,\cdot\cdot,\frac{d}{2}-1;\\ n\in[ 0,|\mathbf{x}|);\end{subarray}}{\text{RoPE}_(n)}= \Big[\ldots,\cos\Big(\mathbb{I}(\hat{\lambda}_{i}, \hat{n})\cdot\frac{n}{\beta^i}\Big),\, \sin\Big(\mathbb{I}(\hat{\lambda}_{i}, \hat{n})\cdot\frac{n}{\beta^i}\Big),\ldots\Big] $}
\\
\text{where } \mathbb{I}(\hat{\lambda}_{i},\hat{n})=
\begin{cases}
1& \text{if } n<\hat{n}\\
{\frac{1}{\lambda}_{i}}& \text{if } n\geq\hat{n}
\end{cases}
\end{gathered}
\end{equation}

In this approach, we define a collection of rescale factors, $\mathbb{I}(\hat{\lambda}_{i},\hat{n})$, that address both types of non-uniformities present. Here, $\hat{\lambda}_i$ refers to the non-uniform scaling for the RoPE dimension indexed by $i$, and $\hat{n}$ represents the token position threshold concerning non-uniformity. The function $\mathbb{I}(\hat{\lambda}_{i},\hat{n})$ adjusts the rotation angle specific to the $i^{th}$ dimension in RoPE by incorporating $\hat\lambda_{i}$ as the scaling parameter and $\hat{n}$ as the threshold for token position. For token positions less than $\hat{n}$, namely the first $\hat{n}-1$ tokens, the original RoPE angle, $\frac{n}{\beta^i}$, remains unchanged and the rescale factor $\hat\lambda_{i}$ is inactive. Once token positions reach or exceed $n \ge \hat{n}$, the rescaling is enabled and $\hat\lambda_{i}$ is employed.

With a desired context window size of $L'$, the aim is to determine the best rescaling coefficients ($\mathbb{I}(\hat{\lambda}_{0},\hat{n})$, $\mathbb{I}(\hat{\lambda}_{1},\hat{n})$, ..., $\mathbb{I}(\hat{\lambda}_{i},\hat{n})$...) across each RoPE dimension from 1 to $d$. This enables the target $\text{LLM}$—utilizing the adjusted $\text{RoPE}$—to reach the lowest possible loss $\mathcal{L}$ (interpreted as perplexity) in next token prediction tasks on input sequences $\mathbf{X}$ of length $L'$.

\subsection{Searching the Non-uniform Position Interpolation}

In order to address the issue presented in Eq.~\ref{eq:problem}, we propose a straightforward but powerful approach that identifies the best RoPE rescale factors. This strategy is designed to take full advantage of the complex, non-uniform characteristics found across multiple dimensions within position embedding.

\textbf{Search space}. Our approach incorporates an expansive search space to capture both types of non-uniformity. Table~\ref{tbl:searchspace} presents the details of this search space. For RoPE embedding, we facilitate the tuning of individualized rescale factors across every dimension. To streamline the construction of the search space, the optimization procedure targets $\lambda_i$ and $\hat{n}$ directly, rather than optimizing $\mathbb{I}(\hat{\lambda}_{i},\hat{n})$, with $\hat{\lambda}_{i}$ defined as $1/\lambda_i$. As depicted in Table~\ref{tbl:searchspace}, the permissible range for $\lambda_i$ spans from a lower bound of 1.0—representing simple extrapolation—to an upper bound of $s\times1.25$, which yields more aggressive interpolation than PI, incremented by 0.01. Here, $s$ denotes the ratio of intended extension for the context window.

The parameter $\hat{n}$ determines how many of the starting token positions maintain their original RoPE embeddings, foregoing any positional interpolation. In practice, we consider $\hat{n}$ values drawn from the set \{0, 1, 2, 4, 8, 12, 16, 20, 24, 28, 32, 64, 128, 256\}. If $\hat{n}=0$, then every token position relies exclusively on the optimized rescale factors.

\textbf{Evolution-based search}. The search space outlined in Table~\ref{tbl:searchspace} encompasses a vast array of positional interpolation options, making thorough exploration highly demanding. For instance, extending to $s=4\times$ yields $400^{128/2}\times14=4\times10^{167}$ potential configurations. As the extension ratio increases, the search space undergoes exponential growth. To manage this complexity, we employ an evolutionary search strategy~\cite{guo2020single} and implement two optimization methods designed to significantly enhance search efficiency. Algorithm~\ref{alg:search} provides an overview of the entire search process.

\textit{Optimized initial population generation}. In place of fully random initialization, our approach inserts the three specific RoPE rescale factors—PI, NTK, and YaRN—directly into the initial population of size $P$ as distinct individuals. The other $P - 3$ slots in the initial population are filled by taking these three factors and applying random mutations to them with a mutation probability $p$.

\textit{Monotonically non-decreasing constraint}. Once the initial population is established, each individual’s LLM perplexity is evaluated. To do this, the appropriate RoPE rescale factors are applied to the target LLM, and the perplexity of the input $\mathbf{X}$ is measured. The highest-performing $k$ individuals are chosen as parents for the evolutionary process. Nonetheless, due to the sheer size of the search space, straightforward mutation and crossover may result in selection of suboptimal candidates, leading to redundant and computationally expensive perplexity assessments. This inefficiency becomes especially pronounced when $L'$ is large, as each calculation of perplexity entails significant inference time.

To overcome this challenge, we enforce a monotonicity constraint such that the sampled RoPE rescaling factors are non-decreasing: $\lambda_i \le \lambda_{i+1}$. Only those RoPE configurations adhering to this property are utilized for perplexity assessment in LLMs, thus streamlining the computational requirements of the search. In particular, we stipulate that $\lambda_i$ must increase monotonically as a function of the RoPE index (i.e., $i=0,\ldots,63$). This constraint is inspired by NTK theory~\cite{jacot2018neural,tancik2020fourier,ntk}, which posits that the lower-frequency dimensions (with smaller $i$) demand reduced interpolation—implying a smaller $\lambda_i$—whereas higher-frequency dimensions permit greater interpolation, justifying larger $\lambda_i$ values.

\begin{algorithm}[t]
	\small
	\caption{The search algorithm}
	\textbf{Input:} target LLM, input samples $\mathbf{X}$, population size $P$, mutation size $N_1$, crossover size $N_2$, maximum iterations $\mathcal{T}$, mutation probability $p$\\

	\begin{algorithmic}[1]
		\label{alg:search}
		\STATE Top-k $\leftarrow \phi$;
		\STATE $\text{P}_0 \leftarrow$ \textit{Initialize\_population\_with\_optimization}($P$, $\mathbf{X}$, $p$);
		\FOR{$i = 1$ \textbf{to} $\mathcal{T}$}
		\STATE \textit{Compute\_perplexity}(LLM, $\text{P}_{i-1}$, $\mathbf{X}$);
		\STATE Top-k $\leftarrow$ \textit{Update\_Topk}(Top-k, $\text{P}_{i-1}$);
		\STATE $\text{P}_{mutation} \leftarrow$ \textit{Mutation\_with\_mono\_constraint}(Top-k, $N_1$, $p$);
		\STATE $\text{P}_{crossover} \leftarrow$ \textit{Crossover\_with\_mono\_constraint}(Top-k, $N_2$);
		\STATE $\text{P}_i \leftarrow \text{P}_{mutation} \cup \text{P}_{crossover} \cup$ Top-k;
		\ENDFOR
		\STATE Output the candidate with the minimum perplexity from Top-k;
	\end{algorithmic}
\end{algorithm}

\textbf{8$\times$ extension without fine-tuning}. Through evolutionary search, our approach uncovers optimal, non-uniform RoPE rescale factors that retain important dimensions and positional encodings, thus reducing information degradation from interpolation. As shown in Fig.\ref{fig:non-ft}, this strategy allows LLaMA2’s context window to expand from 4k to 32k tokens with no fine-tuning required. By comparison, other approaches—including PI, as well as non-uniform NTK and YaRN—see significant increases in perplexity once the context window is extended beyond 2$\times$.

\subsection{Extending LLM Context Window to 2048K}
\label{sec:extendto2048k}

\textbf{Progressive extension to 2048k}. In this section, we present our approach for scaling the context length of pre-trained LLMs beyond the conventional 4k tokens to as much as 2048k. Utilizing non-uniform positional interpolation, we are able to expand the context up to 8$\times$ without requiring any fine-tuning. However, when significantly larger extensions are sought (such as 512$\times$), fine-tuning becomes indispensable. A typical strategy involves determining appropriate RoPE rescaling factors for the target 2048k window followed by fine-tuning. Nonetheless, this process entails substantial computational costs and presents significant training difficulties. Furthermore, as our observations indicate, effectively fine-tuning LLMs for such large extension ratios remains a major challenge (see Appendix).

Fortunately, LongRoPE demonstrates robust performance for both pre-trained and fine-tuned long-context LLMs. Consequently, we propose a progressive and resource-efficient approach capable of reaching the intended 2048k context window, requiring merely 1k fine-tuning steps limited to a 256k training length.

$\Diamond$ \textit{LongRoPE search enables extending pre-trained LLMs to 256k}. For instance, using LLaMA2, we explore expansion to context window sizes of 128k and 256k. In this phase, the extension ratios achieved are 32$\times$ and 64$\times$, respectively.

$\Diamond$ \textit{Fine-tuning to 256k}. In this stage, the pre-trained LLM undergoes fine-tuning to support a 256k context window. Initially, LLaMA2 is fine-tuned for 400 steps with RoPE rescaling set for 128k. Afterward, the RoPE rescaling parameters are updated to those for 256k using the latest checkpoint, followed by a further 600 fine-tuning steps. This approach has shown to be more efficient than commencing fine-tuning directly at the 256k context length.

$\Diamond$ \textit{Scaling fine-tuned extended LLM to 2048k via LongRoPE search}. In the final step, we apply an additional search process on the previously fine-tuned LLM with a 256k token context window. This approach enables a substantial expansion to a context length of 2048k, achieved without additional rounds of fine-tuning. As a result, the context window experiences an overall $512\times$ extension.

\textbf{Restoring performance in shorter context windows}. When expanding the context window to an exceptionally large 2048k tokens, there is a noticeable decline in performance within the initial context range. This phenomenon, previously observed in positional interpolation~\cite{pi}, arises because position embeddings in the original window are compressed into a restricted portion of the embedding space when mapped to higher dimensions. This compression adversely impacts language model performance. The situation is exacerbated by a 512$\times$ scaling, where the embeddings corresponding to tokens in the initial 4k window are densely packed together.

In order to address this issue, we carry out an additional evolution search for the extended LLM, aiming to optimize RoPE rescale factors specifically for shorter context lengths such as 4k and 8k. Because less positional interpolation is necessary at these shorter ranges, we constrain the maximum permissible $\lambda$ during the search process. At inference time, the LLM adapts the relevant RoPE rescale factors in real time.

\section{LongRoPE}

\label{sec:method}
Building on these insights, we propose LongRoPE, which initially incorporates an optimized search strategy designed to leverage both identified non-uniformities. This approach is subsequently utilized to expand the context window of LLMs to support over 2 million tokens.

\subsection{Problem Formulation}

These two forms of non-uniformity generate an expansive solution space and complicate the optimization process. To manage this, we reformulate the multidimensional non-uniform position interpolation optimization as a search problem.

For a LLM targeting a  context window size of $L'$ and lengthy input documents $\mathbf{X}$, where each $\mathbf{x}\in\mathbf{X}$ surpasses $L'$ in token length, we denote the original rotary angle of the $i^{th}$ dimension in RoPE embedding at token position $n$ as  $\frac{n}{\beta^i}$. The optimization problem is then formulated as follows:

\begin{equation}
\begin{gathered}
		\label{eq:problem}

\underset {\mathbf{x}\in\mathbf{X}; \, |\mathbf{x}|\ge L'}{ \operatorname {arg\,min} } \,  \mathcal{L}\, \left( \text{LLM} (\text{RoPE}, \mathbf{X})\right), \;\text{such that} 
\\
\scalebox{0.93}{
$\underset {\begin{subarray}{l} i=0,\cdot\cdot,\frac{d}{2}-1;\\ n\in[ 0,|\mathbf{x}|);\end{subarray}}{\text{RoPE}_(n)}= {\left[\cdot\cdot,cos\left(\mathbb{I}(\hat{\lambda}_{i}, \hat{n})\times\frac{n}{\beta^i}\right), sin\left(\mathbb{I}(\hat{\lambda}_{i}, \hat{n})\times\frac{n}{\beta^i}\right),\cdot\cdot\right] }$}\\
	\text{with} \; \mathbb{I}(\hat{\lambda}_{i},\hat{n})=\begin{cases}
	1&\mbox{\;  $n<\hat{n}$}\\
{\frac{1}{\lambda}_{i}}&\mbox{\; $n\geq\hat{n}$}
\end{cases}
	\end{gathered}
\end{equation}

In this approach, we introduce rescaling factors, $\mathbb{I}(\hat{\lambda}_{i},\hat{n})$, designed to address two distinct types of non-uniformity. Here, $\hat{\lambda_i}$ characterizes the non-uniformity associated with individual RoPE dimensions, while $\hat{n}$ captures non-uniformity in the token positions. Specifically, $\mathbb{I}(\hat{\lambda}_{i},\hat{n})$ serves to adjust the rotational angle for the $i^{th}$ dimension in RoPE, with $\hat{\lambda}_{i}$ representing the scaling parameter and $\hat{n}$ specifying the threshold for token positions. For token positions less than $\hat{n}$, the scaling factor $\hat{\lambda}_{i}$ is not utilized; instead, the conventional RoPE rotary angle $\frac{n}{\beta^i}$ is maintained. Once token positions reach or exceed $n\ge\hat{n}$, the rescaling factor is then introduced.

Suppose a desired context window of length $L'$ is specified. The task is then to determine the optimal set of rescaling factors, represented as $\mathbb{I}(\hat{\lambda}_{0},\hat{n})$, $\mathbb{I}(\hat{\lambda}_{1},\hat{n})$, ..., $\mathbb{I}(\hat{\lambda}_{i},\hat{n})$, for each RoPE dimension from the first to the $d^{th}$. This selection allows the target $\text{LLM}$, utilizing the modified $\text{RoPE}$ embeddings, to reach the lowest possible next-token prediction loss, denoted as $\mathcal{L}$ (corresponding to perplexity), across input sequences $\mathbf{X}$ with token length $L'$.

\subsection{Searching the Non-uniform Position Interpolation}

To address the issue outlined in Eq.~\ref{eq:problem}, we present our straightforward but powerful approach, which involves identifying the best RoPE rescale factors. This method is designed to maximize the benefits of the diverse, non-uniform characteristics present across multiple dimensions in the position embeddings.

\textbf{Search space}. A broad search space has been constructed to incorporate the two types of non-uniformity. Table~\ref{tbl:searchspace} presents an overview of this search space. In particular, a unique rescale factor is permitted for each dimension within the RoPE embedding. To streamline the configuration of the search space, the search is performed over $\lambda_i$ and $\hat{n}$, rather than directly searching for $\mathbb{I}(\hat{\lambda}_{i},\hat{n})$, where  $\hat{\lambda}_{i}=1/\lambda_i$. 
According to Table~\ref{tbl:searchspace}, $\lambda_i$ can take values starting from 1.0 (representing direct extrapolation) up to $s\times1.25$ (corresponding to greater interpolation relative to PI), incremented in steps of 0.01; here, $s$ denotes the ratio by which the context window is extended.

The parameter $\hat{n}$ specifies how many of the earliest token positions retain the standard RoPE embedding, foregoing position interpolation. In practice, we explore $\hat{n}$ values from the set \{0, 1, 2, 4, 8, 12, 16, 20, 24, 28, 32, 64, 128, 256\}. If $\hat{n}=0$, it means that every token position applies the determined rescale factors.

\textbf{Evolution-based search}. The search space detailed in Table~\ref{tbl:searchspace} encompasses a wide range of positional interpolation strategies, which makes thorough exploration highly challenging. For instance, when considering a $s=4\times$ scaling factor, the number of possible configurations rises to $400^{128/2}\times14$ = $4\times10^{167}$. As the extension ratio increases, the search space grows at an exponential rate. To manage this complexity, we adopt an evolution search strategy~\cite{guo2020single} and incorporate two additional optimization methods to significantly improve the efficiency of the search process. The entire search methodology is outlined in Algorithm~\ref{alg:search}.

\textit{Optimized initial population generation}. Rather than using a fully random approach to generate an initial population of $P$ rescale factors, we explicitly include the RoPE rescale factors associated with PI, NTK, and YaRN as three distinct individuals from the outset. The remaining $P - 3$ members of the population are created by applying random mutations to these three rescale factors, each mutation occurring with probability $p$.

\textit{Monotonically non-decreasing constraint}. Upon forming the initial population, we evaluate the perplexity yielded by the LLM for every individual. Concretely, the specified RoPE rescale factors are incorporated into the target LLM, which is then used to assess the perplexity of input $\mathbf{X}$. The top-k scoring individuals are chosen as parents for the evolutionary process. Nonetheless, due to the enormous search space, random mutation and crossover may lead to suboptimal regions, resulting in superfluous perplexity evaluations. This inefficiency becomes more pronounced for large $L'$, as each perplexity computation involves expensive inference.

To overcome this challenge, we enforce a monotonicity condition where the rescaled RoPE factors are non-decreasing: $\lambda_i \leq \lambda_{i+1}$. Only RoPE configurations that adhere to this rule are considered in LLM perplexity evaluations, leading to a notable decrease in search complexity. In essence, $\lambda_i$ must grow monotonically as the RoPE dimension index $i$ progresses from 0 to 63. The rationale for this constraint stems from Neural Tangent Kernel (NTK) theory~\cite{jacot2018neural,tancik2020fourier,ntk}, which posits that dimensions associated with higher frequency (i.e., lower indices) warrant less interpolation (implying a smaller $\lambda_i$), whereas dimensions corresponding to lower frequencies (higher indices) allow for increased interpolation (hence, a larger $\lambda_i$).

\begin{algorithm}[t]
	\small
	\caption{The search algorithm}
	\textbf{Input:} target LLM, input samples $\mathbf{X}$, population size $P$, mutation size $N_1$, crossover size $N_2$, maximum iteration limit $\mathcal{T}$, mutation probability $p$\\

	\begin{algorithmic}[1]
		\label{alg:search}
		\STATE Top-k $\leftarrow \phi$;	
		\STATE $\text{P}_0 \leftarrow$ \textit{Initialize\_population\_with\_optimization} ($P$, $\mathbf{X}$, $p$);
		\FOR{$i=1$ to $\mathcal{T}$}
		\STATE \textit{Compute\_perplexity}(LLM, $\text{P}_{i-1}$, $\mathbf{X}$);
		\STATE Top-k $\leftarrow$ \textit{Update\_Topk}(Top-k, $\text{P}_{i-1}$);
		\STATE $\text{P}_{mutation} \leftarrow$ \textit{Mutation\_with\_mono\_constraint}(Top-k, $N_1$, $p$);
		\STATE $\text{P}_{crossover} \leftarrow$ \textit{Crossover\_with\_mono\_constraint}(Top-k, $N_2$);
		\STATE $\text{P}_i \leftarrow \text{P}_{mutation} \cup \text{P}_{crossover} \cup$ Top-k;
		\ENDFOR
		\STATE Output the individual with the minimum perplexity from Top-k;
	\end{algorithmic}
\end{algorithm}

\textbf{8$\times$ extension without fine-tuning}. Through our evolutionary search, we successfully determine optimal non-uniform RoPE rescale factors that selectively retain crucial dimensions and locations, thereby mitigating information degradation caused by interpolation. As illustrated in Fig.\ref{fig:non-ft}, this approach enables the expansion of LLaMA2’s context window from 4k up to 32k without necessitating fine-tuning. Conversely, previous techniques—including PI, and non-uniform NTK and YaRN—exhibit a sharp increase in perplexity after merely doubling the context size.

\subsection{Extending LLM Context Window to 2048K}
\label{sec:extendto2048k}

\textbf{Progressive extension to 2048k}. Here, we present our approach for scaling the context window of pre-trained LLMs from the standard 4k up to more than 2048k. Our results indicate that non-uniform positional interpolation enables an 8$\times$ increase without the need for any fine-tuning. When even larger expansions (such as 512$\times$) are sought, fine-tuning becomes essential. A commonly considered strategy involves identifying suitable RoPE rescaling parameters for the desired 2048k window before initiating fine-tuning. Nevertheless, this approach is often limited by the significant computational and training costs involved. Additionally, our experience indicates that successfully fine-tuning LLMs for such extensive context expansion presents substantial difficulties (refer to Appendix for more details).

Fortunately, LongRoPE demonstrates effectiveness when applied to both the base model as well as to extended LLMs that have undergone fine-tuning. To this end, we propose a progressive and efficient approach, enabling models to reach the desired 2048k context window using only 1,000 fine-tuning steps, with training sequences not exceeding 256k in length.

$\Diamond$ \textit{LongRoPE search for 256k-length extension in pre-trained LLMs}. Using LLaMA2 as a case study, we perform a search process aiming for context windows of 128k and 256k tokens. At these points, the extension corresponds to a 32$\times$ and 64$\times$ scale-up, respectively.

$\Diamond$ \textit{Fine-tuning to 256k}. Subsequently, we adapt the pre-trained LLM to extend its context window to 256k. In particular, our approach starts by fine-tuning LLaMA2 for 400 steps utilizing RoPE rescaling factors targeted at 128k. After completing this stage, we update the model by substituting the RoPE rescaling factors with those designed for 256k on the obtained checkpoint, followed by a further 600 steps of fine-tuning. Compared to immediate fine-tuning at 256k, this stepwise strategy demonstrates greater efficiency.

$\Diamond$ \textit{Increasing the context length of the fine-tuned extended LLM to 2048k using LongRoPE search}. In the final stage, we apply a second LongRoPE search to the LLM that was previously fine-tuned for sequences up to 256k tokens. This approach allows us to achieve an exceptionally large context window of 2048k tokens, requiring no additional fine-tuning steps. This constitutes a total extension factor of $512\times$.

\textbf{Recovery in shorter context windows}. When the context window is increased to an extensive 2048k length, a decline in model performance is observed within the initial context window. This phenomenon, documented in previous studies on positional interpolation~\cite{pi}, arises because scaling up causes the position embeddings associated with the original context window to be compressed into a smaller subspace of the positional embedding space, which hinders the model’s ability to represent positional information accurately. Especially at an extension ratio of 512$\times$, positional indices mapped from the original 4k context window are densely clustered, exacerbating this compression effect.

To address this issue, we conduct an additional evolution search targeting the expanded LLM, specifically to fine-tune the RoPE rescale parameters for shorter context ranges such as 4k and 8k tokens. For these brief context windows, we lower the upper bound for $\lambda$ during the search, since there is diminished need for positional interpolation. When the LLM is used for inference, it adapts the relevant RoPE rescale values in real-time.

 \section{Experiments}

\subsection{Setup}
\textbf{Evaluation Tasks and models}. The LongRoPE method is integrated with LLaMA2-7B and Mistral-7B, and its effectiveness is assessed across three dimensions: (1) perplexity analysis on long-form texts with extended-context LLMs; (2) the Passkey retrieval task, designed to assess how well the model can locate a straightforward passkey among a substantial amount of extraneous content; and (3) performance on established LLM benchmarks using a standard context window of 4096 tokens.

\textbf{Fine-tuning}. For LLaMA2, we employ a linear learning rate schedule starting at 2e-5, with a global batch size set to 32. The model is initially trained for 400 steps on the Redpajama~\cite{redpajama} dataset, which is segmented into 128k-length chunks, each marked at the beginning and end with BOS and EOS tokens. Subsequently, utilizing the final checkpoint from this phase, we conduct a further 600 training steps in order to extend the context window to 256k tokens. The training for the 128k context window is performed on 8 A100 GPUs using the distributed training framework~\cite{cube}, whereas scaling up to 256k segments necessitates 16 A100 GPUs. For Mistral, we adopt a fixed learning rate of 1e-6 alongside a global batch size of 64. The methodology outlined in YaRN~\cite{yarn} is applied to both the 128k and 256k context models: each undergoes 400 steps of training on Together Computer's Long-Data Collections~\cite{mistraldata}, with sequences of length 16k. Training is carried out using 4 A100 GPUs.

\textbf{Search}. For windows with a target size up to 256k, the following parameters are applied: $P$=64, $N_1$=$N_2$=16, $p$=0.3, $\mathcal{T}$=40, and in each generation, the top-32 candidates are selected for mutation and crossover. The perplexity evaluation utilizes 5 randomly drawn samples from the PG19 validation dataset, ensuring that each sample meets or exceeds the intended context length. When target windows surpass 512k, we reduce the values for population size, mutation, and crossover by half. In this case, perplexity is estimated using 3 randomly chosen samples from the Pile-Books3~\cite{pile} validation split.

\textbf{Baselines}. For achieving a context window of 2048k, models were fine-tuned using context lengths of 128k and 256k, resulting in LongRoPE-2048k (ft=128k) and LongRoPE-2048k (ft=256k) versions for LLaMA2 and Mistral, respectively. The four variants are evaluated alongside leading baselines in large context window extensions, focusing on open-source LLMs that have undergone fine-tuning following positional interpolation with PI, NTK, or YaRN methods. This set of baselines includes Together-32k~\cite{together}, Code LLaMA~\cite{codellama}, LongLoRA-full-FT-100k~\cite{longlora}, as well as YaRN-LLaMA and YaRN-Mistral~\cite{yarn}.

\subsection{Main Results}

\label{sec:mainresult}
\textbf{Long sequence language modeling up to 256k tokens}. To start, we benchmark our approach against leading extended LLMs for sequence lengths up to 256k during evaluation. Our assessment utilizes two different datasets—Proof-pile~\cite{pg19} and the PG19~\cite{pile} test splits—to highlight the generalizability of our methods. We compute perplexity at a range of context lengths, applying a sliding window of 256 tokens. For PG19, evaluations cover the complete test set comprising 100 documents. In the case of Proof-pile, we adhere to YaRN~\cite{yarn}, randomly choosing 10 samples, each spanning a minimum of 128k tokens.

Table~\ref{tbl:proof-pile} and Table~\ref{tbl:pg19} present a perplexity comparison for LLaMA2 and Mistral extended with various interpolation strategies on Proof-pile and PG19, respectively. Two main findings emerge from these results: \textbf{(1)} the perplexity of our extended models consistently decreases as the evaluation length increases from 4k to 256k, demonstrating improved capacity to exploit longer sequences. \textbf{(2)} Despite operating with a context window that is 16$\times$ greater—a condition that typically undermines performance at shorter contexts—our LongRoPE-2048k models still achieve superior results over state-of-the-art baselines up to 256k context length.

\textbf{Long sequence language modeling beyond 2000k}. To assess performance on very large-scale text, we utilize the Books3~\cite{pile} dataset. For more practical evaluation, we randomly choose 20 books, all with lengths over 2048k, and apply a sliding window approach with a size of 256k.

Table~\ref{tbl:books3} demonstrates that LongRoPE effectively expands the context window of LLaMA2-7B and Mistral-7B models up to 2048k tokens, while also attaining perplexity values that are similar to, or better than, the baseline models when evaluated on shorter context sizes between 8k and 128k. Notably, the performance between LLaMA2 at 2048k and Mistral differs significantly: Mistral displays superior results at shorter context windows but its perplexity surpasses 7 after 256k tokens. As expected, LLaMA2's perplexity improves steadily as the window size grows, with only slight increases at the 1024k and 2048k marks. Additionally, for LLaMA2, LongRoPE-2048k exhibits enhanced outcomes when fine-tuned at a window size of 256k compared to 128k, a result attributed to the lower secondary extension ratio (8$\times$ instead of 16$\times$). By comparison, Mistral shows better performance when fine-tuned at a window size of 128k. This difference is mainly because the Mistral fine-tuning for 128k and 256k follows the YaRN protocol, which utilizes a 16k training length, consequently impacting Mistral's capability to extend the context window post-fine-tuning.

\textbf{Passkey retrieval}. Next, we examine how the generation performance depends on the accessible context window size. We employ a synthetic benchmark for passkey retrieval as introduced by~\cite{passkey}. In this benchmark, the model must locate and output a randomly selected five-digit passkey embedded somewhere within an extended document. The precise structure of the prompt can be found in the appendix. For evaluation, we conduct 10 rounds of the retrieval task, with the passkey inserted at uniformly random positions within the tested context window.

As illustrated in Fig.~\ref{fig:passkey}, the retrieval accuracy across different models is compared against baseline performance. Most prior approaches experience a sharp decline in accuracy, falling to zero when the token count exceeds 128k. Conversely, LongRoPE-LLaMA2-2048k (ft=256k) demonstrates remarkable resilience in retrieving passkeys from sequences with millions of tokens, consistently achieving a retrieval accuracy of at least 90\% from 4k to 2048k tokens. Similarly, LongRoPE-Mistral-2048k (ft=128k) sustains perfect accuracy up to 1800k tokens, and only drops to 60\% at 2048k, which is consistent with the observed increase in perplexity at 2048k reported in Table~\ref{tbl:books3}.

\textbf{Standard benchmarks within original context window}. LongRoPE-2048k models are assessed on their original context window utilizing the Hugging Face Open LLM Leaderboard~\cite{huggingfaceboard} under both zero-shot and few-shot conditions. The evaluation includes 25-shot ARC-Challenge~\cite{allenai:arc}, 10-shot HellaSwag~\cite{zellers2019hellaswag}, 5-shot MMLU~\cite{mmlu}, and 0-shot TruthfulQA~\cite{lin2021truthfulqa}.

According to the results in Table~\ref{tbl:commonsense}, our models perform on par with the original benchmark—despite it being tailored for a shorter context window—and surpass the initial Mistral on TruthfulQA by an increase of +0.5\%. While LongRoPE-LLaMA2-2048k, which underwent fine-tuning at 256k, exhibits marginally greater decline in performance, its effectiveness still falls within acceptable bounds across the majority of tasks.

\subsection{Ablation Results}

\textbf{Assessing the efficacy of the second positional interpolation}. Within our progressive extension framework, a secondary application of non-uniform positional interpolation is performed on the fine-tuned extended LLMs through our search algorithm. To empirically demonstrate its advantages, we conduct experiments utilizing the fine-tuned LLaMA2-256k model, subsequently extending it to 512k, 1024k, and 2048k contexts via PI and YaRN techniques. Results presented in Table~\ref{tbl:secondsearch} indicate that our non-uniform positional interpolation maintains perplexity at a steady rate. Conversely, both PI and YaRN exhibit a rapid increase in perplexity as the extension ratio grows.

\textbf{Enhancing recovery at reduced context lengths.} In order to address the decline in performance at lower context lengths, we optimize the RoPE factors in LongRoPE-2048k utilizing our search algorithm. In this process, we set stricter upper limits for the scale factors during the search phase, thereby discouraging excessive interpolation at shorter context ranges such as 4k and 8k. Table~\ref{tbl:shortperformance} presents the perplexity measurements of LongRoPE-LLaMA2-2048k evaluated on Proof-pile at both 4k and 8k, as well as the mean accuracy across standard LLM benchmarks. The findings indicate pronounced improvements in model performance under these short context settings.

\textbf{Investigation of the two types of non-uniformities}. To determine the influence of each non-uniformity component on model performance, we conducted ablation studies. The experimental design included: (i) expanding LLaMA2-7B to shorter context lengths of 16k and 32k using various strategies—PI, optimizing solely the RoPE dimension, and optimizing both non-uniformities; and (ii) scaling our fine-tuned 256k-length LLaMA2 model to 2048k by applying the same approach. Perplexity metrics were measured prior to additional fine-tuning. The results, as seen in Table~\ref{tbl:searchdimension}, demonstrate that addressing non-uniformity within the RoPE dimension substantially decreases perplexity in comparison to PI's linear interpolation. Adjusting for token position non-uniformity also enhances performance on 16k and 32k sequences, but this improvement is not evident at the 2048k length—possibly owing to the vast scale. Simply retaining the earliest tokens without interpolation is ineffective in these settings, which suggests further exploration is warranted for this scenario.

\section{Related Works}

Beyond position interpolation techniques, this section also examines relevant research employing alternative methodologies.

\textbf{Retrieval-based} methods utilize an external memory component for storing extensive historical context and rely on retrieval mechanisms to access pertinent documents during inference~\cite{longllama,wang2023augmenting,borgeaud2022improving}. Such strategies often necessitate deliberate changes to LLM architectural designs. By comparison, our approach is considerably less complex, involving only minimal adjustments to positional embeddings. Furthermore, our method accommodates a broader range of extended context tasks beyond document retrieval, including tasks like long document summarization and few-shot learning.

\textbf{Attention-driven context window expansion}. In addition to interpolating positional embeddings, various studies have enable extending the effective context window for LLMs by altering attention operations within the standard context window size~\cite{han2023lminfinite, streamingllm,parallelwindow}. The central principle here involves employing innovative attention masks to address the challenge of escalating attention computation for additional token positions. Notably, approaches based on attention modification and those leveraging positional interpolation tend to supplement each other.

\textbf{Fine-tuning based approaches} investigate strategies to adapt pre-trained LLMs for extended context capabilities by altering position embeddings. Approaches such as Code LLaMA~\cite{codellama}, LLaMA2 Long~\cite{xiong2023effective}, and ScaledRoPE~\cite{liu2023scaling} employ an enlarged base for RoPE, then fine-tune models on data of the desired context length. Unlike these, our approach maintains adaptability across a wide spectrum of target lengths and can scale to sequences exceeding 2M tokens. With the notable resource intensiveness of fine-tuning for extensive context windows (greater than 128k tokens), recent works like LongLoRA~\cite{longlora} and PoSE~\cite{zhu2023pose} have been introduced to reduce computational demands. Our approach complements these efficient fine-tuning solutions rather than directly overlapping with their methods.

\section{Conclusion}

This paper introduces LongRoPE, a technique that dramatically increases the context length supported by LLMs to a record-breaking 2048k, while preserving their performance in traditional, shorter context ranges. The approach leverages two distinct non-uniform aspects within the RoPE positional embedding, identified via a resource-efficient evolutionary search. This design yields dual advantages: it facilitates robust initialization for subsequent fine-tuning and allows for an 8$\times$ expansion of the context window without necessitating additional fine-tuning. Furthermore, we describe a stepwise extension framework wherein LLMs fine-tuned on 256k-length inputs are progressively scaled to handle a 2048k context window, again avoiding further fine-tuning. Comprehensive experiments demonstrate the efficacy of LongRoPE. We anticipate that LongRoPE-2048k will open the door to a range of novel long-context applications and spark continued innovation in this area.

\section*{Broader Impacts} The research outlined in this paper aims to contribute to progress within the domain of Machine Learning. While our work may have various implications for society, we do not identify any particular impacts that warrant explicit discussion in this section.

	\balance

\end{document}